\documentclass[10pt,a4paper]{amsart}
\usepackage[margin=19mm]{geometry}
\usepackage{amsmath,amssymb,mathtools}
\usepackage[scr=rsfs,cal=euler]{mathalfa}
\usepackage{xfrac}
\usepackage[unicode,hidelinks,bookmarks=false]{hyperref}
\newcommand{\ie}{i.e.\ }
\newcommand{\cf}{cf.\ }
\newcommand{\resp}{respectively\ }
\newcommand{\Subsection}{\subsection}
\providecommand{\nobreakdash}{\nobreak\hbox{-}}
\setlength{\emergencystretch}{2em}
\multlinegap0pt
\usepackage{bm}
\usepackage{bbm}
\usepackage{dsfont}
\usepackage{xspace}
\usepackage{cancel}
\usepackage{mathtools}
\usepackage{amsthm}
\makeatletter
\@ifundefined{theorem}{\newtheorem{theorem}{Theorem}}{}
\@ifundefined{proposition}{\newtheorem{proposition}[theorem]{Proposition}}{}
\makeatother
\newcommand{\CC}{\mathbb{C}}
\title[Equations of Tree Tensor Network Varieties]{Equations of Tree Tensor Network Varieties}
\author{Serkan Ho\c{s}ten, Niharika Chakrabarty Paul, Otto T. P. Schmidt,, Dmitry Skurt}
\date{Source: arXiv:2608.19071v1 (CC BY 4.0)}
\begin{document}
\maketitle

\noindent\emph{Source and licence.} Equations of Tree Tensor Network Varieties. Version arXiv:2608.19071v1 (CC BY 4.0), distributed under the Creative Commons Attribution 4.0 International licence, as recorded in the arXiv licence metadata for this exact version and shown on the abstract page https://arxiv.org/abs/2608.19071v1. Full rights notice: licenses/arxiv-2608.19071-CC-BY-4.0.txt.\par\smallskip

\noindent\textit{Editorial scope.} Counted unit: the Proposition whose source label is \texttt{prop:TT\_is\_TTN} in the article (source0.tex lines 546--564, source label \texttt{prop:TT\_is\_TTN}), delivered together with the article's own complete original proof of it (source0.tex lines 565--637, one \texttt{proof} environment of 73 lines). No mathematics is written, repaired or re-derived: statement and proof are the article's own text line for line. Required context retained uncounted: eq:TTN\_param (lines 372--384); eq:TT-entries (lines 533--543). Every definition, display and label of those lines is retained unchanged. Omitted from this package and not redistributed: 1--371, 385--532, 544--545, 638--2015. Nothing inside a retained excerpt is omitted or altered: every retained body is delivered exactly as the article's lines are written, so no omission receipt of any kind accompanies this package. Citations: the retained excerpts carry no citation command at all, so no bibliography entry of the article is transcribed and no reference is redistributed. Reference adaptation: none was needed and none was made; every reference inside the delivered unit resolves inside it, so no unresolved reference or citation is produced. Printed slips kept literal: no printed word, symbol, hypothesis or number of the counted statement or of its proof was altered. Layout and class: the article's own class is \texttt{article} with packages including amsmath, amsthm, amssymb, color, hyperref, graphicx, caption, mathtools, enumerate, verbatim, mathrsfs, booktabs, tikz, tikz-cd; the wrapper is the lane's pinned \texttt{amsart} editorial wrapper at 10pt on A4, which restates the theorem-like environments the retained text needs and adds the article's own macro definitions that the retained text needs (\texttt{\textbackslash CC}), with extra wrapper packages (bm, bbm, dsfont, xspace, cancel, mathtools). The delivered numbering is the unit's own. Assets: no retained span contains a figure environment, a graphics-inclusion command, a PDF-page inclusion, a table or a TikZ picture; no image file is copied into this package, the package declares no asset dependency and no image file is redistributed with it. Licence: version arXiv:2608.19071v1 of this article is distributed under the Creative Commons Attribution 4.0 International licence, as recorded in the arXiv licence metadata for this exact version and shown on the abstract page https://arxiv.org/abs/2608.19071v1. A full rights notice is delivered at licenses/arxiv-2608.19071-CC-BY-4.0.txt. Characters: every retained excerpt is pure ASCII, so the delivered engine, XeLaTeX with its default Latin Modern text font, reports no missing character.
\begin{align} \label{eq:TTN_param}
\left(\widehat{\Phi}_\rho(\mathbf{A})\right)_{(j_\lambda)_{\lambda\in\mathscr L}} &= \left(
    \widehat{\Phi}_{\mathscr T,\mathbf V,\mathbf r}
    (\mathbf A)
    \right)_{(j_\lambda)_{\lambda\in\mathscr L}} \nonumber\\
    &=
    \sum_{(\alpha_w)_{w\in\mathscr T\setminus\{\rho\}}}
    A^\rho_{(\alpha_u)_{u\in\operatorname{ch}(\rho)}}
    \prod_{v\in\mathscr N\setminus\{\rho\}}
    A^v_{(\alpha_u)_{u\in\operatorname{ch}(v)},\alpha_v}
    \prod_{\lambda\in\mathscr L}
    A^\lambda_{j_\lambda,\alpha_\lambda}.
\end{align}

\begin{equation}\label{eq:TT-entries}
\psi_{j_1,\cdots ,j_N}
=
\sum_{\alpha_1=1}^{r_1}
\cdots
\sum_{\alpha_{N-1}=1}^{r_{N-1}}
A^1_{j_1,\alpha_1}
A^2_{j_2,\alpha_2,\alpha_1}
\cdots
A^N_{j_N,\alpha_{N-1}}.
\end{equation}

\begin{proposition}\label{prop:TT_is_TTN}
Let $\mathbf{d}=(d_1, \ldots, d_N)$ and $\mathbf{r} = (r_0=1, r_1, \ldots, r_{N-1},r_N=1)$ be sequences of positive integers. Then the corresponding TT variety is equal to a TTN variety,
$$ \mathrm{TT}_{\mathbf{d}, \mathbf{r}} = \mathrm{TTN}_{\mathscr{T}, \mathbf{V},\mathbf{r}},$$
where 
\begin{enumerate}
\item $\mathscr{L} = \{\lambda_1,\ldots, \lambda_N\}$ and $\mathscr{N} = \{v_1=\rho, v_2, \ldots, v_N\}$, and the edges in $\mathscr{T}$ are  $\{\lambda_i, v_i\}$ 
for $i=1, \ldots, N$ and $\{v_i, v_{i+1}\}$ for $i=1,\ldots, N-1$,
\item $V_{\lambda_i} = \CC^{d_i}$ for $i=1,\ldots, N$, and
\item $r_{\lambda_i}=d_i$ for $i=1,\ldots,N$, so that
$E_{\lambda_i}=V_{\lambda_i}$, and
\[
r_{v_i}=r_{i-1}
\qquad
\text{for }i=2,\ldots,N.
\]
As before, we use the convention $r_{v_1}=r_\rho=1$.

\end{enumerate}
\end{proposition}

\begin{proof}
Under the rank sequence above, we have $E_{\lambda_i}=V_{\lambda_i}$ for $i \in [N]$ and $E_{v_i}\simeq\mathbb C^{r_{i-1}}$ for $i=2,\ldots,N$, together with $E_{v_1}=E_\rho=\mathbb C$.
First, let $(G^i)_{i=1}^N$ be a collection of ordinary tensor-train
cores, that is, a parametrization realizing a point in $\mathrm{TT}_{\mathbf{d}, \mathbf{r}}$. For every leaf $\lambda_i$, choose the TTN leaf core to be $A^{\lambda_i}
=
\operatorname{id}_{V_{\lambda_i}}.$
Choose the internal TTN core at $v_i$ to be $G^i$, using the
canonical identifications
$$
\begin{aligned}
\mathcal A_\rho
&=
E_{\lambda_1}\otimes E_{v_2}
\simeq
\CC^{d_1}\otimes\mathbb C^{r_1},\\
\mathcal A_{v_i}
&=
E_{\lambda_i}\otimes E_{v_{i+1}}\otimes E_{v_i}^*
\simeq
\CC^{d_i}
\otimes
\mathbb C^{r_i}
\otimes
(\mathbb C^{r_{i-1}})^* \quad \quad i=2,\ldots, N-1,\\
\mathcal A_{v_N}
&=
E_{\lambda_N}\otimes E_{v_N}^*
\simeq
\CC^{d_N}\otimes(\mathbb C^{r_{N-1}})^*.
\end{aligned}
$$
Hence, we realize the tensor-train cores as TTN cores with $G^1 \in \mathcal A_\rho$, $G^i \in \mathcal A_{v_i}$ for $i \in \{2,\ldots, N-1\}$ and $G^N \in \mathcal A_{v_N}$.
Then the TTN contraction \eqref{eq:TTN_param} reduces to the
tensor-train contraction \eqref{eq:TT-entries} and we conclude 
$
\mathrm{TT}_{\mathbf d,\mathbf r}
\subseteq
\mathrm{TTN}_{\mathscr T,\mathbf V,\mathbf r}$.

Conversely, consider TTN cores $(A^v)_{v\in\mathscr{T}}$ that realize a point on $\mathrm{TTN}_{\mathscr T,\mathbf V,\mathbf r}$. Write
$
L^i:=A^{\lambda_i}
\in
V_{\lambda_i}\otimes E_{\lambda_i}^*$
for the leaf core at $\lambda_i$, and write 
$ B^i:=A^{v_i} \in E_{\lambda_i}\otimes E_{v_{i+1}}\otimes E_{v_i}^*
$
for the adjacent internal core. Contract the
$E_{\lambda_i}^*$-factor of $L^i$ with the
$E_{\lambda_i}$-factor of $B^i$ for $i \in [N]$, and define a 
tensor-train core $G^i$ by
\[
G^i_{j_{\lambda_i},\alpha_{v_i},\alpha_{v_{i-1}}}
:=
\sum_{\alpha_{\lambda_i}=1}^{d_i}
L^i_{j_{\lambda_i},\alpha_{\lambda_i}}
B^i_{\alpha_{\lambda_i},\alpha_{v_i},\alpha_{v_{i-1}}} \in
\CC^{d_i}
\otimes
\mathbb C^{r_i}
\otimes
(\mathbb C^{r_{i-1}})^*,
\]
where we use the endpoint conventions $\alpha_0=\alpha_N=1.$
Setting $j_i = j_{\lambda_i}$, $\alpha_i = \alpha_{v_i}$ and substituting these $\mathrm{TT}$ cores into
\eqref{eq:TT-entries} gives exactly the tensor produced by the
original TTN contraction \eqref{eq:TTN_param} for $\mathscr{T}$ defined above. Hence
$
\mathrm{TTN}_{\mathscr T,\mathbf V,\mathbf r}
\subseteq
\mathrm{TT}_{\mathbf d,\mathbf r}$, and
the two varieties are therefore equal.
\end{proof}

\end{document}
